\chapter{Working with the assistant}\label{ch:assistant}
\chapterintro{The assistant panel connects the workspace to a Codex session on your ChatGPT subscription. It builds and
edits the specification through the same tools you use, grounds every region, and asks when a brief leaves a
design-changing fact open. This chapter covers the two autonomy modes, clarification questions with consequence
previews, proposal cards, and what the paper's benchmark says about when this helps.}

\section{What the assistant can do}
Describe a part or a change in words, or paste a whole design brief, and press \keycap{Enter}. Each message starts one
headless Codex session (model \code{gpt-6.1-sol}, reasoning effort high) that has exactly one tool server: the AgenticTO
tool contract. It has no shell, no file access beyond a read-only sandbox, and no other tools, so everything it does is
a logged tool call.

The prompt it receives contains its instructions, the current specification, the item you have selected (so ``this''
and ``it'' refer to your selection), the current rule findings, the latest solve result, and the last turns of the
conversation. Its instructions tell it to:
\begin{itemize}
  \item ground every region with \tool{list_regions} and run \tool{validate_spec} after each change, fixing what it
  reports;
  \item keep the ids that already exist;
  \item when the request is ambiguous in a way that changes the design (which face, how large, what force), not guess but
  call \tool{ask_clarification} once with two to four concrete readings, then apply the reading you choose;
  \item finish with one to three plain sentences: what it changed and anything it assumed.
\end{itemize}

\section{Propose and Act}\label{sec:modes}
\index{Propose mode}\index{Act mode}
The dial in the panel header sets the autonomy (Table~\ref{tab:modes}).

\begin{table}[h]
\caption{The two autonomy modes.}\label{tab:modes}
\begin{fixtable}{@{}>{\raggedright\arraybackslash}p{0.13\linewidth}>{\raggedright\arraybackslash}X>{\raggedright\arraybackslash}X@{}}
Mode & What the assistant may do & How you see it \\ \midrule
\gui{Propose} (default) & Works on a private copy of your specification: a stand-alone session seeded with it, with the
\code{clarify} tool profile. It can edit, ground, validate and ask; it cannot solve. & When it finishes, every
difference from your specification becomes a card. Nothing changes until you accept. \\
\gui{Act} & Attached to the live session with the \code{full} profile. Its edits apply at once. Its instructions allow
the sketch and preview rungs to check its work and leave refine and release to you. & Edits appear immediately in the
tree, the viewport and the timeline, tagged \gui{AGENT}. \gui{Undo this run's edits} reverts the whole run in one
step. \\
\end{fixtable}
\end{table}

\begin{notebox}[An instruction, not a lock]
In Act mode the restriction to sketch and preview is part of the assistant's instructions. The \code{full} profile it
uses does contain \tool{solve} for every rung. Even if an agent did release a design, the certificate would still come
from the solver, and you would see the solve in the timeline.
\end{notebox}

\section{Clarification questions with consequence previews}\label{sec:clarify}
\index{clarification}\index{consequence preview}
When a brief leaves something open that changes the design, the assistant asks instead of guessing. Each possible
reading is a concrete option, either a complete support, load or keep-in/keep-out item, or parameter values such as a
volume fraction. The workspace solves each reading on the sketch rung, all starting from the same latest design so that
their compliances compare fairly, and shows it as a chip with a thumbnail of the design it leads to and its compliance
(Figures~\ref{fig:clarify} and~\ref{fig:chips}). Pick a chip, or type your own answer under \gui{None of these}.

\begin{figure}[t]
\wideshot{20_clarify.png}
\caption{Building a part from the ambiguous brief b44, the assistant asks which axis is vertical and which way the
torque turns. Each reading is shown with the design it produces. In Propose mode your own specification is untouched
until you accept, so the viewport still shows the empty starting box.}
\label{fig:clarify}
\end{figure}

\begin{figure}[tb]
\begin{minipage}[b]{0.4\linewidth}\shot[\linewidth]{21_clarify_chips.png}\end{minipage}\hfill
\begin{minipage}[b]{0.55\linewidth}
\caption{The clarification chips, each with its sketch design and compliance. Picking a chip answers the question
with that reading.}
\label{fig:chips}
\end{minipage}
\end{figure}

The question, its options, the previews and your answer are written to the session log, so the chosen interpretation
becomes part of the record, and a replay reuses the recorded answer without asking again.

\begin{tipbox}[Take your time]
A question stays open for ten minutes. The assistant is told to keep waiting for your answer. If its session has ended
by the time you answer, your answer starts a follow-up run that applies it, so the answer is not lost.
\end{tipbox}

\section{Cards and the ghost overlay}\label{sec:cards}
\index{cards}
In Propose mode the result is a set of cards, one per change: \emph{add}, \emph{change} or \emph{remove} for an item,
\emph{set} for a field such as the volume fraction. Hover a card to see its ghost in the viewport
(Figure~\ref{fig:cards}). \gui{Accept} or \gui{Reject} each card, or use \gui{Accept all} and \gui{Reject all}. Accepted
cards are applied as one undoable step, logged as yours (actor \code{user}), with the cards and the agent run recorded
beside it.

\begin{figure}[t]
\wideshot{22_cards.png}
\caption{Proposed changes as cards. Hovering a card draws its ghost in the viewport.}
\label{fig:cards}
\end{figure}

\begin{figure}[t]
\wideshot{23_accepted.png}
\caption{After \gui{Accept all}: the specification is built and valid, and the ladder is ready.}
\label{fig:accepted}
\end{figure}

\section{An Act run}
In Figure~\ref{fig:act} the assistant in Act mode lowers the volume fraction and runs a sketch to check the design. The
steps drawer lists every tool call. A failed call shows its error, which the assistant then fixes. The footer gives the
wall time and the token use. \gui{Undo this run's edits} reverts the run as a whole.

\begin{figure}[p]
\wideshot[150mm]{24_act.png}
\caption{An Act run: six tool calls (list the regions, set the volume fraction, validate, solve the sketch, wait, read the field
statistics), the 16,000-element sketch in the viewport, and \gui{AGENT} records in the timeline. The run took 14~s.}
\label{fig:act}
\vspace{8mm}
\begin{minipage}[b]{0.30\linewidth}\shot[\linewidth]{25_act_panel.png}\end{minipage}\hfill
\begin{minipage}[b]{0.65\linewidth}
\caption{The assistant panel after an Act run: the instruction, the steps drawer with the six numbered tool calls and
their arguments, the link that undoes the run, the footer with wall time and token use, and the reply.}
\label{fig:actpanel}
\end{minipage}
\end{figure}

\section{When the assistant helps, and what it costs}\label{sec:x2}
The paper measured how often agents recover the intended specification from a brief, on a suite of 50 three-dimensional
briefs: 40 to solve, 7 that need a clarification, and 3 that must be refused because they are ill-posed as stated.

\begin{evidencebox}[Specification grounding on 50 briefs, primary provider, first repeat]
Right behaviour means a specification equivalent to the intended or an admissible reading, or a refusal where the
brief must be refused \claim{C2.1, C2.3}.

\smallskip
{\centering\sffamily\footnotesize
\begin{tabular}{@{}lrr@{}}
\toprule
Condition & Right behaviour & Silent wrong \\ \midrule
T: text only, JSON from the brief and schema & 37 / 50 & 1 \\
G: grounding tools, no questions allowed & 44 / 50 & 6 \\
G+C: grounding and clarification & 50 / 50 & 0 \\
G+C+V: G+C followed by a three-view picture check & 50 / 50 & 0 \\
\bottomrule
\end{tabular}\par}

\smallskip
T to G corrected 8 tasks and broke 1 ($p \approx 0.039$); G to G+C corrected 6 and broke none ($p \approx 0.031$)
\claim{C2.2}. A deterministic rule parser without a language model got 9 right.
\end{evidencebox}

Three lessons for use:
\begin{enumerate}
  \item \textbf{Grounding fixes representation errors.} Text-only output often failed to encode valid geometry (for
  example zero-length cylinders for washers). With the tools these errors largely disappear.
  \item \textbf{Clarification fixes intent errors, at a price.} The G condition was not allowed to ask and guessed; its
  6 silent errors all passed the rules. Allowed to ask, the agent fixed all six, but it asked 59 questions, 45 of them on
  fully stated briefs \claim{C2.5}. Expect to confirm details that the brief already states.
  \item \textbf{The benefit depends on the agent choosing to ask.} A second provider (Claude Code, \code{claude-opus-5-5},
  effort high) with the same tools asked 7 questions instead of 59 and reached 46 of 50 under G+C, keeping 4 silent
  wrong specifications; two of them were unit slips. The difference in right outcomes was not significant on this sample
  ($p \approx 0.12$) \claim{C2.6}.
\end{enumerate}

\begin{warnbox}[Read the cards, not the summary]
An agent can explain its assumption fluently and still return a specification that differs from what you meant. The
cards, the tree and the viewport show what the specification actually says. Accept a card only when you would have
written that item yourself.
\end{warnbox}

\begin{notebox}[What the benchmark does and does not show]
In the benchmark a scripted oracle answered the questions from the stored reference reading; it ignored the
consequence previews. The results therefore show that asking and using the answer works under a cooperative oracle,
not that the previews improve a person's choice. The reference readings were adjudicated by two authors, the second blind to condition and provider, and no
participant study has been run.
\end{notebox}
